% ============================================================================
% COURS 03 — MODÉLISATION DES TRANSMETTEURS — PARTIE B
% Source : 4-ModelisationTransmetteurs2022.docx
% ============================================================================

\section{Transmetteurs avec transformation de mouvement}
\label{sec:trans-avec}

\subsection{Pignon-crémaillère}

Le pignon est guidé en rotation et la crémaillère en translation. Le
roulement sans glissement entre le cercle primitif et la ligne de
référence impose :

\begin{center}\TLPignonCremaillere\end{center}

\[
\boxed{v=\pm R\omega=\pm\frac{mZ}{2}\,\omega.}
\]
Le signe dépend des orientations choisies. Le dispositif est réversible
et transforme une rotation continue en translation continue.

\begin{TLThreeApp}{Borne rétractable}
Une borne est actionnée par un moteur, un réducteur de rapport $1/60$ et
un pignon-crémaillère de rayon primitif $R$.
\[
v_{\text{borne}}=R\,\omega_{\text{pignon}}
=\frac{R}{60}\,\omega_{\text{moteur}}.
\]
Déterminer la vitesse du moteur nécessaire pour une vitesse de levée
imposée.
\end{TLThreeApp}

\subsection{Cylindre roulant sans glisser sur un plan}

Un cylindre de rayon $R$ roulant sans glisser se comporte
cinématiquement comme un pignon sur une crémaillère :
\[
\boxed{|v_O|=R|\omega|.}
\]
Le point de contact avec le plan a une vitesse nulle. Cette relation ne
doit pas être confondue avec la vitesse d'un point quelconque de la
périphérie.

\subsection{Courroie ou chaîne utilisée comme organe linéaire}

Lorsque deux poulies de même diamètre déplacent un brin rectiligne, la
translation du brin est reliée à la rotation d'une poulie par :
\[
\boxed{v=R\omega.}
\]
La solution est réversible, économique pour les grands entraxes et
courante dans les convoyeurs, imprimantes, portiques et machines-outils.

\subsection{Vis-écrou}

La liaison hélicoïdale couple une rotation et une translation suivant le
même axe.

\begin{center}\TLVisEcrou\end{center}

Pour un pas $p$ correspondant à l'avance par tour :
\[
\boxed{x=\frac{p}{2\pi}\theta,}
\qquad
\boxed{v=\frac{p}{2\pi}\omega.}
\]
Le signe dépend du sens de l'hélice et du solide immobilisé. Une vis à
billes réduit les frottements, améliore le rendement et peut rendre le
système réversible ; une vis trapézoïdale peut être irréversible.

\begin{TLThreeApp}{Pilote automatique de bateau}
Une vis de pas $p=4\ \mathrm{mm}$ est entraînée par une machine à courant
continu à travers un réducteur de rapport $r$. La vitesse de la tige est :
\[
v_{4/0}=\frac{p}{2\pi}\,r\,\omega_{15/0}.
\]
Identifier tous les étages de la chaîne et déterminer le coefficient
global en $\mathrm{m\,rad^{-1}}$.
\end{TLThreeApp}

\section{Train d'engrenages épicycloïdal}
\label{sec:trans-epi}

Un train épicycloïdal permet d'obtenir un rapport important dans un
faible encombrement, avec des arbres d'entrée et de sortie souvent
coaxiaux.

\begin{center}\TLTrainEpicycloidal\end{center}

\subsection{Constituants}

\begin{center}
\small
\begin{tabularx}{\linewidth}{>{\bfseries}p{3cm}X}
\toprule
Constituant & Définition \\
\midrule
Porte-satellite & pièce en rotation d'axe fixe par rapport au bâti, portant les axes des satellites \\
Satellite & roue dentée dont l'axe est fixe par rapport au porte-satellite mais mobile par rapport au bâti \\
Planétaire & roue ou couronne d'axe fixe par rapport au bâti, engrenant avec les satellites \\
\bottomrule
\end{tabularx}
\end{center}

Plusieurs satellites identiques ne modifient pas la loi cinématique ; ils
répartissent les efforts et augmentent la puissance transmissible.

\subsection{Raison de base}

On imagine le porte-satellite bloqué. Le train devient un train simple à
axes fixes. La \textbf{raison de base} est alors :
\[
\boxed{\lambda=
\left.\frac{\omega_A}{\omega_B}\right|_{\omega_{PS}=0}.}
\]
Elle se calcule avec la formule des trains simples en tenant compte des
contacts intérieurs et extérieurs.

\subsection{Formule de Willis}

La relation générale entre les deux planétaires $A$, $B$ et le
porte-satellite $PS$ est :
\[
\boxed{\frac{\omega_A-\omega_{PS}}
{\omega_B-\omega_{PS}}=\lambda.}
\]
Une forme développée particulièrement pratique est :
\[
\boxed{\omega_A-\lambda\omega_B+(\lambda-1)\omega_{PS}=0.}
\]
La somme des coefficients de cette équation est nulle, propriété utile
pour vérifier l'écriture.

\begin{TLThreeMethod}{Étudier un train épicycloïdal}
\begin{enumerate}
  \item Repérer les deux planétaires, le ou les satellites et le porte-satellite.
  \item Immobiliser mentalement le porte-satellite et calculer $\lambda$ comme pour un train simple.
  \item Écrire la formule de Willis.
  \item Introduire la condition de fonctionnement : planétaire bloqué, porte-satellite bloqué ou vitesse imposée.
  \item Résoudre pour obtenir le rapport demandé.
  \item Vérifier le signe, l'ordre de grandeur et les cas limites.
\end{enumerate}
\end{TLThreeMethod}

\subsection{Configurations usuelles}

Un train simple comporte trois arbres cinématiques possibles : les deux
planétaires et le porte-satellite. Pour obtenir une transmission à un
degré de liberté, il faut imposer deux vitesses parmi les trois, souvent
en bloquant l'un des arbres.

\begin{center}
\small
\begin{tabularx}{\linewidth}{>{\bfseries}p{3.6cm}X}
\toprule
Élément bloqué & Conséquence \\
\midrule
Planétaire intérieur & entrée et sortie possibles entre l'autre planétaire et le porte-satellite \\
Planétaire extérieur / couronne & architecture fréquente de réducteur coaxial \\
Porte-satellite & retour à un train simple à axes fixes \\
Aucun élément bloqué & différentiel ou sommateur de vitesses à deux entrées indépendantes \\
\bottomrule
\end{tabularx}
\end{center}

\begin{TLThreeApp}{Trois configurations}
Pour un train de raison de base $\lambda$, déterminer successivement :
\begin{enumerate}
  \item $\omega_{PS}/\omega_A$ lorsque $B$ est bloqué ;
  \item $\omega_B/\omega_A$ lorsque $PS$ est bloqué ;
  \item $\omega_{PS}/\omega_B$ lorsque $A$ est bloqué.
\end{enumerate}
Les résultats doivent être déduits de la même équation de Willis, sans
reconstruire la cinématique à chaque fois.
\end{TLThreeApp}

\subsection{Différentiel}

Dans un différentiel, aucune des trois branches n'est nécessairement
bloquée. Une vitesse d'entrée est distribuée entre deux sorties dont les
vitesses peuvent être différentes. L'équation de Willis impose leur
relation moyenne. Le couple est réparti selon l'architecture et les
frottements internes.

\section{Dispositifs complémentaires}
\label{sec:trans-complementaires}

\begin{center}
\small
\begin{tabularx}{\linewidth}{>{\bfseries}p{3cm}p{3.5cm}p{2.2cm}X}
\toprule
Dispositif & Transformation & Réversibilité & Usages \\
\midrule
Bielle-manivelle & rotation continue $\leftrightarrow$ translation alternative & parfois & moteurs, compresseurs, pompes \\
Pignon-crémaillère & rotation continue $\leftrightarrow$ translation continue & toujours & direction, portes, axes linéaires \\
Vis-écrou & rotation continue $\leftrightarrow$ translation continue & selon technologie & vérins électriques, machines-outils \\
Croix de Malte & rotation continue $\rightarrow$ rotation intermittente & non & indexage de plateaux \\
Excentrique & rotation continue $\rightarrow$ translation alternative & généralement non & pompes, ablocage, taille-haie \\
Came axiale & rotation continue $\rightarrow$ translation programmée & non & pompes et automatismes \\
\bottomrule
\end{tabularx}
\end{center}

\begin{TLThreeDef}{Bielle-manivelle et approximation de grande bielle}
Pour une manivelle d'excentration $e$ et une bielle de longueur $L$, la
translation est non sinusoïdale. Lorsque $L\gg e$, typiquement
$L>5e$, le mouvement se rapproche d'une loi sinusoïdale et certaines
expressions peuvent être simplifiées.
\end{TLThreeDef}

\begin{TLThreeDef}{Croix de Malte}
L'ergot moteur entraîne la roue à rainures pendant une fraction du tour,
puis la roue reste immobile pendant le reste de la rotation. L'angle
entre rainures fixe le nombre de positions d'indexage.
\end{TLThreeDef}

\section{Synthèse des lois usuelles}
\label{sec:trans-synthese}

\begin{center}
\small
\renewcommand{\arraystretch}{1.25}
\begin{tabularx}{\linewidth}{>{\bfseries}p{3cm}p{4.8cm}X}
\toprule
Transmetteur & Loi cinématique & Point de vigilance \\
\midrule
Roues de friction & $\omega_2/\omega_1=-D_1/D_2$ & adhérence et effort presseur \\
Engrenage extérieur & $\omega_2/\omega_1=-Z_1/Z_2$ & même module \\
Engrenage intérieur & $\omega_2/\omega_1=+Z_1/Z_2$ & même sens de rotation \\
Train simple & $(-1)^n\prod Z_m/\prod Z_M$ & roues folles et contacts intérieurs \\
Roue-vis sans fin & $|r|=Z_v/Z_r$ & rendement et irréversibilité \\
Poulies / chaîne & $\omega_2/\omega_1=D_1/D_2$ & glissement possible pour courroie \\
Pignon-crémaillère & $v=R\omega$ & signe selon paramétrage \\
Vis-écrou & $v=p\omega/(2\pi)$ & pas, sens de l'hélice, réversibilité \\
Train épicycloïdal & $(\omega_A-\omega_{PS})/(\omega_B-\omega_{PS})=\lambda$ & identifier correctement les constituants \\
\bottomrule
\end{tabularx}
\end{center}

\section{Sources de la partie transmetteurs}
\label{sec:trans-sources}

Cette partie reprend le cours original consacré aux transmetteurs, des
ressources techniques publiques et des activités de collègues de
l'UPSTI. Les exercices conservent les systèmes et les données du
document Word.
